\subsection{Finding 1: Cross-Indicator Inconsistency}
\label{sec:finding1}

Tethered to a single-indicator paradigm, existing visual token pruning methods rest upon an assumption that has largely escaped scrutiny: a single importance criterion suffices to separate indispensable tokens from expendable ones.
Our investigation reveals that the indicators employed by existing methods naturally coalesce into two distinct families: \textit{saliency}, which quantifies individual token importance (VisionZip~\cite{visionzip}, CDPruner~\cite{cdpruner}), and \textit{diversity}, which promotes collective representativeness of the retained set (DivPrune~\cite{divprune}, SCOPE~\cite{scope}).

As illustrated in Fig.~\ref{fig:indicator_vis}, under an identical compression ratio, these two indicator families yield markedly divergent token selections.
VisionZip, guided by ViT CLS attention, preferentially retains foreground tokens of high semantic saliency---a selection that remains invariant to the user's query.
CDPruner augments this saliency signal with text--visual cosine similarity, redirecting retention toward query-relevant regions while its conditional DPP kernel enforces geometric spread among the selected tokens.
By contrast, DivPrune pursues a fundamentally different objective: maximizing the minimum pairwise distance in the feature space, it scatters selections broadly across the image at the expense of fine-grained foreground detail.
SCOPE mediates between these two poles, coupling CLS-based saliency with a submodular coverage term that curbs token clustering within narrow semantic neighborhoods.

To be specific, methods sharing the same indicator family exhibit moderate concordance---for instance, VisionZip and SCOPE, both grounded in CLS-based saliency, achieve an overlap of 0.64 at $K{=}128$.
Cross-family overlap, however, deteriorates precipitously: at $K{=}32$, it plummets to a mere 0.19--0.27.
Even under the most permissive retention budget ($K{=}128$, retaining 22\% of all visual tokens), the average pairwise overlap remains below 0.45---implying that \emph{more than half of the tokens deemed ``most important'' by one indicator family are judged expendable by the other}.
This pronounced inconsistency reveals that XXX.

This divergence arises because the four methods optimize along fundamentally \textit{orthogonal indicator dimensions}:

\begin{itemize}[leftmargin=*]
  \item \textbf{Dimension~A---Saliency} evaluates individual token importance and further decomposes into two sub-dimensions with distinct behaviors:
  \begin{itemize}
    \item \textbf{A1 (Vision-Intrinsic):} \textit{``Is this token intrinsically important in the image, regardless of the question?''}---relying solely on vision-encoder signals (\eg, CLS attention in VisionZip and SCOPE), producing a fixed ranking for a given image.
    \item \textbf{A2 (Task Relevance):} \textit{``Is this token relevant to the user's current query?''}---requiring text--visual interaction (\eg, CDPruner computes cosine similarity between each visual token and the instruction embedding), yielding query-dependent rankings that differ across prompts.
  \end{itemize}
  \item \textbf{Dimension~B---Representativeness} evaluates whether the selected subset collectively represents the original token set, and further decomposes along two complementary spaces:
  \begin{itemize}
    \item \textbf{B1 (Semantic Diversity):}
          \textit{``Are the selected tokens semantically non-redundant in the embedding space?''}---enforced via set-theoretic objectives that penalize feature-space duplication (\eg, DivPrune's MMDP maximizes the minimum pairwise distance; SCOPE's submodular coverage maximizes the sum of nearest-neighbor similarities; CDPruner's conditional DPP maximizes the volume spanned by the selected feature vectors).
    \item \textbf{B2 (Spatial Coverage):}
          \textit{``Does the selected set preserve at least one representative token from every spatial region of the image?''}---enforced via geometric partitioning of the token grid (\eg, Nuwa~\cite{nuwa} divides the grid into $M{\times}M$ non-overlapping blocks and selects top-$n$ tokens per block; HoloV~\cite{holov} allocates per-crop budgets proportional to holistic scores).
          Notably, B1 and B2 are not redundant: a semantically diverse subset may still concentrate in a single spatial quadrant, while a spatially uniform selection may retain redundant tokens from texturally homogeneous regions.
  \end{itemize}
\end{itemize}

Crucially, these dimensions are not merely orthogonal---their failure modes are \emph{complementary}.
Optimizing saliency alone (Dimension~A) causes the retained tokens to cluster around a few high-attention foreground regions, leaving the semantic space severely under-covered; SCOPE's analysis quantifies this: under pure CLS-attention selection on LLaVA-1.5, the $\theta$-coverage of the retained set falls \emph{below that of random sampling} across all threshold values~\cite{scope}.
Optimizing semantic diversity alone (B1) yields the opposite pathology: DivPrune distributes tokens broadly in the embedding space but is agnostic to which regions are task-relevant, retaining background noise at the expense of fine-grained foreground details.
Furthermore, even a semantically diverse subset (B1) may remain spatially skewed---tokens from texturally varied but co-located patches can achieve high pairwise distance in feature space while leaving entire image quadrants unrepresented, a systematic weakness for OCR, spatial reasoning, and referring expression comprehension~\cite{nuwa}.

These complementary failure modes reveal that no single indicator dimension suffices: saliency without representativeness discards context; representativeness without saliency discards detail; and semantic diversity without spatial awareness discards regions.
This analysis motivates the design of a unified pruning framework that jointly optimizes across all four sub-dimensions (A1, A2, B1, B2), as we present in Sec.~\ref{sec:method}.
